% !TeX program = xelatex
\documentclass[11pt,a4paper]{article}
\usepackage[margin=23mm]{geometry}
\usepackage{fontspec}
\setmainfont{Latin Modern Roman}
\usepackage{amsmath,amssymb,mathtools}
\usepackage{xcolor,enumitem,booktabs,longtable,array,xurl,hyperref}
\definecolor{muted}{HTML}{53636C}
\hypersetup{colorlinks=true,urlcolor=blue,linkcolor=blue}
\setlength{\parindent}{0pt}
\setlength{\parskip}{5pt}
\newcommand{\R}{\mathbb R}
\newcommand{\E}{\mathbb E}
\newcommand{\N}{\mathbb N}
\newcommand{\ind}{\mathbf 1}
\newcommand{\Var}{\operatorname{Var}}
\newcommand{\Cov}{\operatorname{Cov}}
\newcommand{\supp}{\operatorname{supp}}
\DeclareMathOperator*{\argmin}{arg\,min}
\DeclareMathOperator*{\argmax}{arg\,max}
\newcommand{\entrymeta}[3]{{\sffamily\small\color{muted}\textbf{\texttt{#1}}\quad|\quad #2\par #3\par}}
\newcommand{\Problemref}[1]{\texttt{#1}}
\newcommand{\JELref}[2]{\texttt{#2} (JEL \texttt{#1})}
\begin{document}
\section*{Simple lending contracts preserving group-finance benefits}
\textbf{Problem ID:} \texttt{G21-198}\quad \textbf{JEL:} \texttt{G21}\par
% BEGIN ECONOMICS STATEMENT
\label{problem:OP-0647}


\label{audited:ECON-086}

{\sffamily\small\color{muted}\textbf{ECON-086} | R | Source-backed agenda; exact formalization is editorial\par}

\subsection*{Definitions and formal target}

This is a parameterized research question. For each fully specified economic model \(M\), characterize the frontier conditional on its supplied primitives and conventions. An instance specifies a fixed baseline borrower population \(I\), type domains for \(\theta_i\), repayment-action domains for \(r_i\), project and support technologies, horizon, timing, information, enforcement, and a measurable contract space \(\mathcal A_M\). Contract \(a\) specifies liability, meeting requirements and repayment rules. Feasible group memberships, including nonmembership, are distinguished from the social network \(G\) through which borrowers exchange information or support; membership need not be encoded by that network. The instance specifies whether no lending, solo borrowing and empty groups are available.

The model supplies a formation and repayment game or other behavioral specification, its equilibrium concept, and an equilibrium selection rule, including any randomization. Under each contract for which the selected equilibrium exists, this determines membership, peer actions, repayment and the possibly random network \(G_M(a)\). Formation is endogenous when these objects respond to the contract through that specification. Define \(\mathcal A_M^{\mathrm v}\subseteq\mathcal A_M\) to contain precisely contracts with a well-defined selected equilibrium satisfying the model's resource, enforcement and behavioral constraints and its individual voluntary-participation requirements. The instance states when participation is assessed, the information on which it is conditional, and the available refusal and exit options. Coalition-stability requirements apply only if explicitly included in the model. A contract with no equilibrium satisfying these requirements is outside this domain; an unspecified equilibrium concept or selection leaves the instance incomplete.

In the display below, \(G=G_M(a)\) and \(U(a,G)\) and \(\Pi(a,G)\) denote the resulting random borrower-welfare and net lender-payoff accounts, with other equilibrium arguments suppressed. The instance supplies individual utilities, outside payoffs, aggregation and welfare weights over the fixed baseline population, including the outcomes of excluded borrowers and nonparticipants. It specifies lender funding, administration, meeting and enforcement costs and all support transfers or real support costs. Each real cost is counted once. The law underlying \(\mathbb E_M\) includes types, correlated project risks, endogenous transitions and selection randomization as applicable. Payoff accounts specify their horizon and discounting, are integrable, and have not already averaged over randomness that is integrated again in the display.

Complexity \(K:\mathcal A_M\to\mathbb R\) is a supplied, real-valued assessment index with a declared scale and budget domain. Its assessment procedure states whether it measures written rules or also depends on borrower characteristics, network outcomes or implementation; any such dependence is evaluated by a predeclared rule to obtain \(K(a)\). The index is defined independently of the desired welfare ranking. If its underlying burden also entails a utility or resource cost, that cost enters the payoff accounts once; imposing the complexity bound does not itself subtract a cost.

For each supplied instance \(M\) and declared complexity budget \(\overline K\), characterize the frontier
\[
 \max_{a\in\mathcal A_M^{\mathrm v}:\,K(a)\leq\overline K}\mathbb E_M[U(a,G_M(a))]
 \quad\text{s.t.}\quad\mathbb E_M[\Pi(a,G_M(a))]\geq0.
\]
Here the frontier comprises the optimal borrower-welfare value as a function of the budget and the feasible contracts attaining it, or approaching its supremum when it is not attained. Report infeasible budgets, finite nonattained suprema and unbounded welfare separately. No lending or an empty group ensures feasibility only if it is an available outcome satisfying the stated constraints.

Including endogenous formation and the value of social support, determine which simplifications preserve valuable features of group finance. For this comparison an instance additionally supplies a reference contract, a relation specifying which contracts count as its simplifications, and a preservation criterion on their induced outcomes. It identifies the valued support channels and their welfare or benefit functionals, and states whether preservation concerns aggregate welfare, specified components, or a specified counterfactual comparison. A counterfactual comparison specifies what is changed and what is held fixed, with formation and subsequent behavior evaluated under each compared contract. Characterize the feasible simplifications satisfying that supplied criterion and their relationship to the frontier. These comparison inputs are required choices; no particular benchmark, tolerance, support mechanism or universal contract ranking is supplied by the citations. The requested characterization is structural and conditional on the declared inputs, with no additional algorithmic or empirical-identification requirement. Without those inputs the question remains a research agenda rather than an evaluated optimization or preservation claim.

\subsection*{Clarifications and required assumptions}

Predeclare a complexity metric \(K\), such as the number of state-contingent rules, measured comprehension burden or implementation cost; these are not interchangeable. Define each borrower's utility, information, outside option and repayment behavior, liability and lender costs. Group formation, peer actions and repayment require a stated equilibrium concept and selection, especially with strategic interactions. Include borrower participation and lender viability; positive average welfare does not guarantee voluntary membership. A fixed network \(G\) differs from a network formed under the contract.

\subsection*{Variants and changes of scope}

\begin{itemize}
\item Joint liability, meetings, information-sharing and mutual insurance are separate group-finance components.
\item Fixed groups versus endogenous formation change selection and incentives.
\item Simplicity can be a constraint or valued outcome, but its measure must be independently defined.
\end{itemize}

\subsection*{Audit conclusion and source relationship}

[S1] and [S2] support simple contracts preserving valuable group/social-capital features. This frontier is editorial and overlaps mechanism design and game theory.

\subsection*{References and exact source roles}

{\footnotesize\raggedright
\par\noindent\textbf{[S1]} Jing Cai, Muhammad Meki, Simon Quinn, Erica Field, Cynthia Kinnan, Jonathan Morduch, Jonathan de Quidt, and Farah Said (2025). \emph{Microfinance}. VoxDevLit 3(3), January 2025.\par \textit{Verified locator / source role:} Primary publisher page checked for review identity, authors, version and background. The specific published agenda is corroborated in [S2]; the exact formal target is editorial.\par \url{https://voxdev.org/voxdevlit/microfinance}\par\smallskip
\par\noindent\textbf{[S2]} Oliver Hanney / VoxDev (living compilation). \emph{Evidence gaps: Key unanswered questions in development economics}. Accessed 8 October 2026. The page currently covers 20 topics; the ``16-topics URL is a historical slug.\par \textit{Verified locator / source role:} Topic heading: Microfinance. Supports the broad research agenda; this is not a verbatim quotation or an attribution of the displayed mathematical target.\par \url{https://voxdev.org/topic/evidence-gaps-key-unanswered-questions-across-16-topics-development-economics}\par\smallskip
}

\subsection*{Common conventions: Additional-source status and mathematical conventions}

\label{audited:conventions}

Of the 100 supplied ECON entries, 65 distinct problems appear here; 32 close
matches refine earlier entries, and three duplicate questions map to existing
entries. Appendix~\ref{app:audited-crosswalk} lists every destination.
The attachment's P/R/S/U distinctions and source audits are retained. Its
precise mathematical formalizations are editorial unless the individual audit
says otherwise. Candidate subsidy targets ECON-004--006 retain their U flags;
the resolved approval-core question ECON-007 updates OP-0202 above.
The following supplied conventions govern both this part and the attributed
ECON refinements elsewhere in the catalogue, subject to local restrictions.

\textbf{Parameterized questions.} A problem family lists inputs that must be fixed before an instance is evaluated: population, horizon, policies, information, data law, model class and welfare weights. These are required choices, not quantities the citations are assumed to specify. Undefined applications should not be treated as identified estimands.

\textbf{Indices and integrability.} Sets of agents, firms and regions are finite unless a population measure is explicitly used. \(i,j\) index units, \(r\) regions, \(t\) dates and \(h\) horizons. \(\mathbb E\) and \(\Pr\) refer to the declared population and law; \(\mathbf1\{A\}\) is an indicator. Every displayed expectation and discounted sum needs existence in the stated domain. Infinite horizons use a discount in \((0,1)\) and a convergence condition; finite horizons may use discount one.

\textbf{Optimization and existence.} A displayed \(\max\), \(\min\), or \(\arg\max\) is conditional on attainment. For finite feasible menus this is immediate; general spaces need continuity and compactness/coercivity or another existence argument. Without attainment, the target is the supremum/infimum and arbitrarily near-optimal feasible choices. \(\inf\varnothing=+\infty\). Empty feasibility and an unidentified target are different issues.

\textbf{Potential outcomes.} \(Y_i(z)\) is the outcome under a specified intervention, not the observed conditional mean among people choosing \(z\). In binary comparisons \(\Delta Y_i=Y_i(1)-Y_i(0)\). Specify assignment, treatment versions, compliance, information and observation horizon. Interference is allowed under a full policy vector; compressed exposure notation needs a justified mapping. No interference, consistency and overlap are substantive assumptions, not automatic consequences of notation.

\textbf{Populations and observation.} Fix baseline eligibility and population weights. Conditioning on post-policy employment, survival, relocation or program uptake can change the population being compared. Death is not ordinary loss to follow-up; outcomes truncated by death need a composite convention, joint survival target, or explicitly defined survivor stratum. Fertility-changing interventions need a population functional rather than an unexplained fixed descendant cohort. Measurement and missingness models must be supplied.

\textbf{Identification.} A structural model \(m\in\mathcal M\) implies observable law \(P_m\) and target \(\theta(m)\). Identification means
\[
 P_{m_1}=P_{m_2}\ \Longrightarrow\ \theta(m_1)=\theta(m_2)
 \quad\text{for all }m_1,m_2\in\mathcal M .
\]
Without identification, the target is
\[
 \Theta(P;\mathcal M)=\{\theta(m):m\in\mathcal M,\ P_m=P\}.
\]
Writing \(\theta(P)\) before establishing identification can hide incompatible latent models. Confidence sets additionally need a sampling law, confidence level, asymptotic sequence and uniformity class. Point identification, coverage and informative precision are separate properties.

\textbf{Mediation.} Total effects of randomized assignment are distinct from cross-world natural indirect effects. Belief/effort-channel effects require a meaningful mediator intervention and additional measurement, exclusion or mediation assumptions. Randomization of the initial policy alone does not supply those assumptions. Report bounds or interventional alternatives when the stronger target is unsupported.

\textbf{Equilibrium.} A counterfactual \(W(z)\), \(p(z)\), or \(\mu_t^z\) requires individual optimization, feasibility, government/resource budgets, market clearing and state-transition laws. State equilibrium existence and selection, or report ranges over compatible equilibria. Initial-state integration includes subsequent endogenous transitions; current-state integration must use the policy-dependent distribution. Reduced-form invariance after policy is not assumed.

\textbf{Welfare accounts.} Scores, utility, health, dollars and ecological quantities require explicit conversion before addition. Fees, taxes, subsidies, wages and extortion payments are transfers between parties; distributional weights can make them welfare relevant, but their face value is not automatically a resource loss. State ownership, geographic boundaries, foreign-income weights and fiscal finance. Do not add profits to household welfare already including dividends, or count land capitalization and the underlying benefit twice. A benefit-cost expression must distinguish gross benefits from net welfare and subtract every real cost exactly once.

\textbf{Derivatives and risk.} Log derivatives require positive arguments; local responses require admissible perturbations and specified equilibrium branches. Differentiation under expectations needs regularity/domination, and boundaries may allow only one-sided derivatives. Risk uses a specified law; ambiguity uses a declared nonempty law set on a common history space. Adaptive policies must be nonanticipative. A robust criterion is an editorial choice unless attributed explicitly.

\textbf{Scope overlaps.} Allocation, collective choice, contracts, household bargaining and coordinated policy overlap economics with optimization and game theory. They are retained because they appeared in the original catalogue; no claim of a strictly game-theory-free selection is made.
% END ECONOMICS STATEMENT
\end{document}
